\documentclass[11pt,a4paper]{article}
\usepackage[utf8]{inputenc}
\usepackage[T1]{fontenc}
\usepackage[margin=2.2cm]{geometry}
\usepackage{lmodern}
\usepackage{microtype}
\usepackage{amsmath,amssymb}
\usepackage[table]{xcolor}
\usepackage{booktabs}
\usepackage{tabularx}
\usepackage{float}
\usepackage{enumitem}
\usepackage[most]{tcolorbox}
\usepackage{fancyhdr}
\usepackage{needspace}
\usepackage{titlesec}
\usepackage{tikz}
\usetikzlibrary{arrows.meta,positioning,calc,shapes.geometric,fit}
\usepackage[colorlinks=true,linkcolor=blue!60!black,urlcolor=blue!60!black]{hyperref}

\setlist[itemize]{nosep,leftmargin=*,topsep=2pt,partopsep=0pt}
\setlist[enumerate]{nosep,leftmargin=*,topsep=2pt,partopsep=0pt}
\widowpenalty=10000
\clubpenalty=10000
\setlength{\parskip}{0.55em}
\setlength{\parindent}{0pt}

\definecolor{cblue}{HTML}{2B6CB0}
\definecolor{cgreen}{HTML}{2F855A}
\definecolor{cred}{HTML}{C53030}
\definecolor{corange}{HTML}{DD6B20}
\definecolor{cgray}{HTML}{718096}
\definecolor{cpurple}{HTML}{6B46C1}

\titleformat{\section}{\Large\bfseries\color{cblue}}{\thesection}{0.6em}{}
\titleformat{\subsection}{\large\bfseries\color{cblue!80!black}}{\thesubsection}{0.6em}{}

\pagestyle{fancy}
\fancyhf{}
\lhead{\small Paper 12 explained simply}
\rhead{\small Chatterjee \& Govindu, arXiv 1505.01936}
\cfoot{\small\thepage}
\renewcommand{\headrulewidth}{0.3pt}

\newtcolorbox{idea}[1][]{colback=cblue!6,colframe=cblue!70,fonttitle=\bfseries,title={#1},boxrule=0.6pt,arc=2pt,left=6pt,right=6pt,top=4pt,bottom=4pt}
\newtcolorbox{toy}[1][]{colback=corange!7,colframe=corange!80,fonttitle=\bfseries,title={#1},boxrule=0.6pt,arc=2pt,left=6pt,right=6pt,top=4pt,bottom=4pt}
\newtcolorbox{warn}[1][]{colback=cred!6,colframe=cred!75,fonttitle=\bfseries,title={#1},boxrule=0.6pt,arc=2pt,left=6pt,right=6pt,top=4pt,bottom=4pt}
\newtcolorbox{eqbox}[1][]{colback=cgray!8,colframe=cgray!70,fonttitle=\bfseries,title={#1},boxrule=0.6pt,arc=2pt,left=6pt,right=6pt,top=4pt,bottom=4pt}
\newtcolorbox{ours}[1][]{colback=cgreen!7,colframe=cgreen!75,fonttitle=\bfseries,title={#1},boxrule=0.6pt,arc=2pt,left=6pt,right=6pt,top=4pt,bottom=4pt}

\begin{document}
\setcounter{tocdepth}{1}

% ------------------------------------------------------------ TITLE
\begin{center}
{\LARGE\bfseries\color{cblue} Paper 12, explained from zero}\\[0.6em]
{\large \textit{Noise in Structured-Light Stereo Depth Cameras: Modeling and its Applications}}\\[0.5em]
Avishek Chatterjee and Venu Madhav Govindu\\
Department of Electrical Engineering, Indian Institute of Science, Bengaluru\\
arXiv 1505.01936 (8 May 2015), 31 pages\\[0.4em]
\textbf{Reading order \#3 of the ball-estimation chapter}
\end{center}

\begin{idea}[What this guide is]
This is a slow, step-by-step explanation of one paper, written for someone who is \emph{not}
used to academic papers, but it keeps the academic vocabulary so you can later read the
original. Statements that come from the paper are marked \textbf{(paper)}. Blue boxes are the
main ideas. Orange boxes (\textbf{Toy example}) use numbers that I computed myself (checked by
computer) or invented; they are \emph{not} from the paper unless marked. Red boxes are my own
critical comments. Green boxes say what this means for \emph{our} MuJoCo ball chapter.

I read the full text of the paper (31 pages) with \texttt{pdftotext}. The paper has three
``applications'' in Section~3 (denoising, scan merging, plane extraction) that are about 3D
reconstruction; I explain them briefly, because only their \emph{idea} (weight by uncertainty)
matters for us.
\end{idea}

\subsection*{The whole paper in six sentences}
A structured-light depth camera such as the Microsoft Kinect measures depth by triangulation: an
infrared projector shines a known dot pattern, an infrared camera sees where each dot lands, and
the shift of a dot (the \emph{disparity}) tells how far the surface is. The disparity can only be
measured with a fixed precision, which the authors argue is the main noise source and does
\emph{not} depend on distance. Because depth is the \emph{inverse} of disparity, a fixed error in
disparity becomes a depth error that grows with the \textbf{square of the distance}:
$\sigma_Z\propto Z^2$. They check this on a real Kinect by showing that the \emph{spacing} between
the distinct depth values the camera can output grows with a slope of $1.967$ on a log-log plot
(theory says $2$). They then use the model to build a depth-aware smoothing filter, a better way
of merging scans, and a plane detector. The paper does \textbf{not} say anything about moving
objects, dropouts, delays or small balls, and it gives no final numeric noise level.

\tableofcontents
\newpage

% ============================================================ 1
\section{Why we read this paper, and what to take from it}

\begin{center}
\begin{tabularx}{\textwidth}{@{}c p{4.6cm} X@{}}
\toprule
 & Question for our chapter & What this paper gives us \\
\midrule
1 & How does depth noise depend on distance? & The law $\sigma_Z = k\,Z^2$, with a physical derivation (Sec.~\ref{sec:derivation}). \\
2 & What is the constant $k$? & Not given directly. We can \emph{derive} a baseline from the Kinect numbers in the paper (Sec.~\ref{sec:k}); I label it as my derivation. \\
3 & Does depth also have a coarse ``step'' size? & Yes: depth comes in discrete steps that grow like $Z^2$ (Sec.~\ref{sec:quant}). \\
4 & How should a filter weight measurements at different distances? & Inverse variance, so weight $\propto 1/Z^4$ (Sec.~\ref{sec:apps}). \\
5 & Is this noise important for \emph{our} camera? & Surprisingly little, for a downward-looking camera and a planar ball (Sec.~\ref{sec:ours}). \\
\bottomrule
\end{tabularx}
\end{center}

\begin{warn}[Four cautions before we start]
(1)~The camera is the \textbf{Kinect (version 1)}, a structured-light device from about 2010--2012.
Our simulated camera is a generic overhead depth camera. (2)~The paper studies \textbf{static
scenes} (floor, statues, rooms); nothing moves. (3)~It never states the size of the disparity noise
in pixels, only that disparity is estimated with $1/8$-pixel steps. (4)~Its numbers are about a
sensor used \emph{indoors at 0.5--3~m}; the paper itself says such cameras ``do not work in the
presence of moderate to strong sunlight''.
\end{warn}

% ============================================================ 2
\section{The story in one page}

Close one eye and hold a finger in front of your face. Open that eye and close the other: the
finger seems to jump sideways relative to the background. The \emph{closer} the finger, the
\emph{bigger} the jump. Your brain uses this jump to feel depth. A structured-light camera does
the same with a projector as one ``eye'' and an infrared camera as the other, 75~mm apart.

For each little patch of the image it measures the jump (the \textbf{disparity} $D$, in pixels).
A big jump means near, a small jump means far. The depth is
\[
Z=\frac{fB}{D},
\]
where $f$ is the focal length in pixels and $B$ is the distance between the two ``eyes'' (the
\textbf{baseline}). Now the trap: measuring the jump has a small, fixed error, say a fraction of
a pixel. For a nearby surface the jump is large (e.g.\ 70 pixels), so a fraction of a pixel is
nothing. For a far surface the jump is tiny (e.g.\ 10 pixels), and the same fraction of a pixel is
a large share of it. Because $Z$ comes from \emph{dividing} by $D$, small errors on a small $D$
turn into large errors in $Z$. This is the entire idea of the paper.

\begin{figure}[H]
\centering
\begin{tikzpicture}[scale=1.0, every node/.style={font=\small}]
\draw[->] (0,0) -- (8,0) node[right]{depth $Z$};
\draw[->] (0,0) -- (0,4.2) node[above]{depth error $\sigma_Z$};
\draw[very thick, cblue, domain=0.3:7.5, samples=60] plot (\x, {0.065*\x*\x});
\node[cblue] at (5.6,3.9) {$\sigma_Z = k\,Z^2$};
\draw[dashed, cgray] (2,0) -- (2,0.26); \draw[dashed, cgray] (6,0) -- (6,2.34);
\node[below] at (2,0) {near}; \node[below] at (6,0) {far};
\node[right, text width=3cm, align=left, font=\footnotesize] at (6.3,1.2) {$3\times$ farther\\ $=9\times$ noisier};
\end{tikzpicture}
\caption{The main result: depth error grows with the square of the distance.}
\end{figure}

\begin{idea}[One sentence to remember]
\emph{A depth camera that triangulates has a constant error in disparity but an error in depth
that grows as the square of the distance, so the same camera is precise up close and coarse far
away.}
\end{idea}

% ============================================================ 3
\section{Glossary}

\begin{center}
\begin{tabularx}{\textwidth}{@{}l X@{}}
\toprule
Term & Meaning \\
\midrule
Structured light & Projecting a known pattern (here infrared dots) and looking at how it deforms on the scene. \\
Stereo / triangulation & Finding depth from two viewpoints separated by a known distance. \\
Disparity $D$ & The shift, in pixels, of a feature between the two viewpoints. Large when near, small when far. \\
Baseline $B$ & The distance between the two viewpoints (projector and camera). Kinect: $75$~mm \textbf{(paper)}. \\
Focal length $f$ & The camera's zoom, expressed in pixels. Kinect IR camera: $587$ px \textbf{(paper)}. \\
Depth map & An image in which every pixel holds a distance instead of a colour. \\
Quantisation & Rounding a continuous value to a grid of allowed values. Like measuring with a ruler that only shows whole millimetres. \\
Sub-pixel resolution & Measuring disparity in fractions of a pixel. Kinect: $1/8$ pixel \textbf{(paper)}. \\
Standard deviation $\sigma$ & The typical size of a random error. \\
Additive noise & Noise that is just added to the true value: $D = D_0 + n$. \\
Bilateral filter & A smoothing filter that averages neighbouring pixels but gives less weight to neighbours with very different values, so it keeps edges. \\
TSDF, volumetric merging & A way to fuse several depth scans into one 3D model. Not needed for us. \\
MLE & Maximum likelihood estimate: the value that makes the observed data most probable. For Gaussian noise with different variances it is the \emph{inverse-variance weighted average}. \\
Principal point $(c_x,c_y)$ & The pixel where the optical axis hits the image; the image centre. The paper writes $(u,v)$. \\
\bottomrule
\end{tabularx}
\end{center}

% ============================================================ 4
\section{The paper, section by section}

\subsection{Introduction (paper, Sec.~1)}
The authors survey earlier work on Kinect depth noise. Two earlier studies had already noticed
\emph{empirically} that the noise grows quadratically with depth; two others derived a noise model
from the working principle but did not use it. The authors say their contribution is to derive the
model, validate it, \textbf{and use it} in three applications.

\subsection{The noise model: derivation, in tiny steps (paper, Sec.~2.1)}
\label{sec:derivation}
\textbf{Step 1: the camera equation.} Treat the projector as a second camera. The depth is
\[
Z=\frac{fB}{D}. \tag{paper: $D=fB/Z$}
\]
\textbf{Step 2: the noise assumption.} The estimated disparity is the true disparity plus an
additive error: $D = D_0 + n$, where $n$ has a \textbf{fixed} standard deviation $\sigma_D$, the
same for all distances. The authors say the dominant part of $n$ is \emph{quantisation noise}
(rounding of the disparity to $1/8$ pixel).

\textbf{Step 3: put it into the depth formula.}
\[
Z=\frac{fB}{D_0+n}=\frac{fB}{D_0\,(1+n/D_0)}\ \approx\ \frac{fB}{D_0}\Bigl(1-\frac{n}{D_0}\Bigr)
= Z_0-\frac{fB}{D_0^{2}}\,n .
\]
The approximation is the rule $1/(1+\varepsilon)\approx 1-\varepsilon$ for small $\varepsilon$, which
holds because the noise $n$ is small compared with the disparity $D_0$.

\textbf{Step 4: replace $D_0$.} From $D_0=fB/Z_0$ we get $fB/D_0^{2}=Z_0^{2}/(fB)$. So
\[
Z\approx Z_0-\frac{Z_0^{2}}{fB}\,n
\qquad\Longrightarrow\qquad
\boxed{\ \sigma_Z=\frac{Z_0^{2}}{fB}\,\sigma_D\ } \tag{paper Eq.~1}
\]
The depth error is the disparity error multiplied by $Z_0^2/(fB)$. The authors derive the same
thing by differentiating $D=fB/Z$: $\partial Z/\partial D=-Z^{2}/(fB)$ (paper Eq.~2).

\begin{eqbox}[In plain words]
A one-pixel mistake in disparity costs \emph{more and more depth} the farther the surface is. The
conversion factor is $Z^2/(fB)$ millimetres per pixel of disparity.
\end{eqbox}

\textbf{The paper's example (paper).} With $B=75$~mm and $f=587$~px:
\[
\frac{\partial Z}{\partial D}\Big|_{Z=600\,\mathrm{mm}}=-\frac{600^2}{75\cdot587}=-8.2\ \text{mm/px},\qquad
\frac{\partial Z}{\partial D}\Big|_{Z=1500\,\mathrm{mm}}=-\frac{1500^2}{75\cdot587}=-51.1\ \text{mm/px}.
\]
So one pixel of disparity error is $8.2$~mm at $2$ feet and $51.1$~mm at $5$ feet. I recomputed
both numbers and they are right.

\subsection{Empirical validation (paper, Sec.~2.2)}
\label{sec:quant}
How do you check that the noise grows like $Z^2$ on a real Kinect? The authors scan a flat floor,
which contains a continuous range of depths, and list all the \emph{distinct} depth values the
camera outputs, in increasing order. If depth resolution were uniform, the gaps between neighbouring
distinct values would be constant. Instead the gaps grow with distance. They plot
$\log(\Delta Z)$ against $\log(Z)$ (where $\Delta Z$ is the gap between consecutive distinct values)
and fit a line: the slope is \textbf{1.967} (paper), nearly $2$, as the theory predicts.

They also find out how finely the Kinect measures disparity: sorting the disparities, there are at
most \textbf{8} distinct values between two integers, so the \textbf{sub-pixel resolution is
$1/8$ pixel} (paper). Near the camera some values coincide because the output depth is also rounded
to $1$~mm. They then \emph{simulate} the pipeline in a few lines of Matlab (paper, Algorithm~1)
and reproduce the log-log plot. In Python the same four lines read:
\begin{center}\small\texttt{z = b*f / (np.round(b*f/np.arange(500,3001)*8)/8);\ \ z = np.unique(np.round(z));}\\
\texttt{plot(np.log(z[:-1]), np.log(np.diff(z)))}\quad with $b=75$, $f=587$.\end{center}

\begin{toy}[Toy example: the depth ``staircase'' of such a camera]
Disparity is rounded to steps of $\delta=1/8$ pixel. A step of $\delta$ in disparity is a depth step
of $\Delta Z = Z^2\,\delta/(fB)$. With $fB=587\cdot0.075=44.0$~px$\cdot$m:
\begin{center}\small
\begin{tabular}{@{}rrr@{}}
\toprule
distance $Z$ & depth step $\Delta Z$ (1/8 px) & noise if only rounding: $\sigma_Z=\Delta Z/\sqrt{12}$ \\
\midrule
1.5 m & 6.4 mm & 1.8 mm \\
3.0 m & 25.6 mm & 7.4 mm \\
4.0 m & 45.4 mm & 13.1 mm \\
\bottomrule
\end{tabular}
\end{center}
(The rule $\sigma=\text{step}/\sqrt{12}$ is the standard deviation of a uniformly distributed
rounding error.) \textbf{This is my derivation, not a number in the paper.} It says that at $3$~m
this camera can only report depth in steps of about $2.6$~cm, close to the diameter of our ball
($6$~cm), while the random error from rounding alone is about $7$~mm.
\end{toy}

\subsection{Three applications (paper, Sec.~3)}
\label{sec:apps}
All three use the same idea: \emph{do not treat all depth values as equally trustworthy}.
\begin{enumerate}
\item \textbf{Adaptive bilateral filter (Sec.~3.2).} A bilateral filter smooths a depth image but
keeps edges; it has a threshold $\sigma_d$ that decides when two neighbouring depths ``belong
together''. With a fixed threshold, near regions are smoothed correctly but far regions are not
smoothed enough (or near regions are over-smoothed). Setting $\sigma_d(p)=k\,Z(p)^2$ with a
constant $k$ makes it work at all distances (paper Eq.~6). The paper shows this on pictures; it
gives no numeric error table.
\item \textbf{Weighted volumetric scan merging (Sec.~3.3).} When several scans of the same
surface are fused, a scan taken from far away should count less. For independent Gaussian errors
with different standard deviations $\sigma_i$, the best average weights each measurement by
$1/\sigma_i^2$. Since $\sigma_i\propto Z_i^2$, the weight is
\[
w_i=\frac{1}{Z_i^{4}} \tag{paper Eq.~13}
\]
(a constant factor cancels). A statue scanned from $75$~cm and $1.5$~m shows the benefit: equal
weights wash out the details of the close scan.
\item \textbf{Plane extraction (Sec.~3.4--3.5).} A test ``is this point within $T$ of the plane?''
needs a different $T$ for near and far points. The trick: \emph{disparity} has constant noise, and
for a plane the disparity is an affine function of the pixel coordinates (paper Eq.~15), so
working in disparity space allows one fixed threshold.
\end{enumerate}

\begin{eqbox}[The one formula we take from these applications]
For independent measurements with standard deviations $\sigma_i$, the best combination is
\[
\hat\mu=\frac{\sum_i x_i/\sigma_i^{2}}{\sum_i 1/\sigma_i^{2}},\qquad \sigma_i\propto Z_i^2
\ \Rightarrow\ \text{weights}\ \propto Z_i^{-4}.
\]
A Kalman filter does exactly this when its measurement-noise covariance $R$ grows with the
variance of the measurement. So a filter fed depth from different distances should use
$R\propto Z^{4}$.
\end{eqbox}

\subsection{Conclusion (paper, Sec.~4)}
The authors restate the model, its validation, and its three uses. There is no discussion of
limitations or of other camera types.

% ============================================================ 5
\section{Numbers worth remembering}

\begin{center}
\small
\begin{tabularx}{\textwidth}{@{}X l@{}}
\toprule
Item & Value \\
\midrule
Baseline $B$ (Kinect) & $75$~mm \textbf{(paper)} \\
IR focal length $f$ & $587$ px, field of view about $57^\circ$ \textbf{(paper)} \\
Depth sensitivity at $600$~mm / $1500$~mm & $8.2$ / $51.1$ mm per pixel of disparity \textbf{(paper)} \\
Sub-pixel resolution of disparity & $1/8$ px \textbf{(paper)} \\
Depth output precision & $1$~mm \textbf{(paper)} \\
Slope of $\log\Delta Z$ against $\log Z$ & $1.967$, theory $2$ \textbf{(paper)} \\
Weight for scan merging & $w=1/Z^4$ \textbf{(paper)} \\
Noise constant $k$ (or $\sigma_D$ in pixels) & \textbf{not given} \\
\bottomrule
\end{tabularx}
\end{center}

% ============================================================ 6
\section{Critical reading (my comments)}

\begin{warn}[1. What the experiment really validates]
The figure with slope $1.967$ validates that the \emph{spacing between output depth values} grows
as $Z^2$. That is the quantisation step, not the random error. The statement that the standard
deviation $\sigma_Z$ grows as $Z^2$ then follows from the model (a constant disparity noise),
and the authors never measure $\sigma_Z$ against distance directly (for instance by scanning a
plane many times and computing the spread). It is a reasonable argument, but it is a model
inference.
\end{warn}

\begin{warn}[2. ``Quantisation dominates'' is asserted, not shown]
The paper says the ``most significant source of disparity error is quantization noise'' and lists
others (sensor noise, ambient infrared, surface reflectance). No experiment separates them. If
other noise sources were larger, $k$ would be larger, although the $Z^2$ shape would stay.
\end{warn}

\begin{warn}[3. One sensor, one era, static scenes]
The law comes from triangulation, so it holds for other structured-light or active-stereo cameras.
It does \emph{not} apply to time-of-flight cameras, whose noise behaves differently, and the paper
tests only the Kinect. It also says nothing about moving objects, frame-to-frame behaviour,
\emph{holes} (pixels with no depth) except for small ones near occlusions, or reflective and
dark surfaces.
\end{warn}

\begin{warn}[4. Noise is assumed independent per pixel]
The model treats every depth value as having its own independent error. The paper does not say
whether neighbouring pixels share errors. For a ball covering about fifty pixels, independent
errors could be averaged down by $\sqrt{50}$, while shared errors could not.
\end{warn}

% ============================================================ 7
\section{What this means for our chapter}
\label{sec:ours}

\begin{ours}[1. The model we use]
Depth noise in the sensor stage: $Z_{\text{measured}} = Z + \mathcal N(0,\sigma_Z^2)$ with
$\sigma_Z=k\,Z^2$. This is exactly the structured-light law of this paper.
\end{ours}

\begin{ours}[2. A baseline value for $k$ (my derivation, to be confirmed)]
\label{sec:k}
From the Kinect numbers ($B=75$~mm, $f=587$~px) and rounding noise only ($\sigma_D=0.125/\sqrt{12}=
0.036$~px): $k=\sigma_D/(fB)=8.2\times10^{-4}$~m$^{-1}$, which gives $\sigma_Z=1.8$~mm at $1.5$~m
and $7.4$~mm at $3$~m. If we use our own focal length ($416$~px) with the Kinect baseline, $k$
becomes $1.2\times10^{-3}$ and $\sigma_Z=10.4$~mm at $3$~m. Because the true $k$ of a real camera is
larger than the rounding-only value, I propose to \textbf{sweep $k$} over $\times1$, $\times5$,
$\times20$ of the baseline instead of trusting one number.
\end{ours}

\begin{toy}[Toy example, measured on our own camera: how much does depth noise move the ball?]
In our overhead camera the ball position is $x=(u-c_x)\,Z/f$. A depth error $\delta Z$ therefore
moves the planar position by $\delta x=\delta Z\cdot (u-c_x)/f = \delta Z\cdot x/Z$: \emph{zero at
the image centre and about $\delta Z/3$ at the arena edge} ($x=1$~m, $Z=3$~m). I tested this with
our \texttt{DepthCamera} (3~m, $60^{\circ}$, 150 ball positions inside the arena, noise added to
the depth read at the ball's centroid; scratch test, not in the project):
\begin{center}\small
\begin{tabular}{@{}lrr@{}}
\toprule
condition & mean $xy$ error & max $xy$ error \\
\midrule
ideal depth (no noise) & 1.7 mm & 4.0 mm \\
depth noise at the baseline ($\sigma_Z=7.4$~mm, $\times1$) & 2.2 mm & 8.9 mm \\
depth noise $\times5$ ($\sigma_Z=37$~mm) & 7.8 mm & 35 mm \\
depth noise $\times20$ ($\sigma_Z=148$~mm) & 31 mm & 134 mm \\
\textbf{no depth noise, centroid jitter $0.3$~px} & 3.2 mm & 8.6 mm \\
\bottomrule
\end{tabular}
\end{center}
At the paper's rounding-level noise, depth error hardly matters for the \emph{planar} position
(1.7 to 2.2~mm). A mere $0.3$-pixel jitter of the detected centroid hurts more (3.2~mm). Depth noise
becomes important only if it is $5$ times larger or more.
\end{toy}

\begin{warn}[What that toy does and does not show]
It shows that, for a \emph{downward-looking camera and a planar ball}, the dominant imperfections
are not depth noise but (a)~pixel-level detection error, (b)~lost frames, (c)~delay and
(d)~frame rate. Depth noise would matter more if the camera were tilted, mounted on the arm, or if
we detected the ball from depth instead of colour. It also does not show the staircase effect:
a coarse $2.6$~cm depth step shifts all pixels of the ball surface by the same amount, which
produces a position \emph{bias} that depends on where the ball is, not random noise. (In my test
the quantised case happened to give a smaller error than the ideal case; that is a coincidence of
where the steps fall, not an improvement.)
\end{warn}

\begin{ours}[3. What goes into the sensor stage]
Proposed switches for Step~1D (each independently on or off):
\begin{enumerate}
\item \textbf{Depth noise} $\sigma_Z=kZ^2$ with $k$ as above (sweep $\times1,5,20$).
\item \textbf{Depth quantisation} with step $Z^2/(8fB)$ (optional).
\item \textbf{Centroid jitter} in pixels (to be justified from paper 18; placeholder $0.3$~px).
\item \textbf{Dropout}, flagged and silent hold (decided in the paper~11 deep dive).
\item \textbf{Delay} (placeholder $40$~ms, swept) and \textbf{frame rate} ($60$~Hz).
\end{enumerate}
\end{ours}

\begin{ours}[4. For the estimators]
Use a measurement covariance $R$ that grows with the depth variance, $R\propto Z^4$ in the depth
direction. For an overhead camera $Z$ is almost constant (about $2.94$ to $3$~m), so $R$ is nearly
constant; it would matter for an arm-mounted camera in the later step.
\end{ours}

\begin{eqbox}[Open decision for you]
Which imperfections should we put into the \emph{main} experiments? My proposal: depth noise
(baseline $k$, sweeping $\times1,5,20$), centroid jitter, flagged dropout, delay and frame rate,
with depth quantisation off by default. If you prefer to simulate a specific real camera (for
instance the Kinect with $B=75$~mm), say so and I will use its parameters.
\end{eqbox}

% ============================================================ 8
\section{Test yourself}

\begin{enumerate}
\item Why does a constant error in disparity become an error in depth that grows with the square
of the distance? Derive $\sigma_Z=Z^2\sigma_D/(fB)$ from $Z=fB/D$.
\item With $B=75$~mm and $f=587$~px, what is the depth error caused by one pixel of disparity error
at $Z=1$~m? (Answer: $1000^2/(75\cdot587)=22.7$~mm.)
\item What does the log-log slope of $1.967$ validate, and what does it \emph{not} validate?
\item Why should a Kalman filter that gets depth measurements from different distances use
$R\propto Z^4$?
\item For a camera 3~m above the floor and a ball at $x=1.5$~m from the optical axis, how does a
$1$~cm depth error change the measured $x$? (Answer: $\delta x=\delta Z\cdot x/Z=5$~mm.)
\item Name three imperfections of our simulated camera that this paper does \emph{not} cover.
\item Why is depth noise less important for a downward-looking camera than for a tilted one?
\end{enumerate}

\appendix
\section{Map of the paper}
\begin{center}
\small
\begin{tabularx}{\textwidth}{@{}l X@{}}
\toprule
Paper part & Content and where it is in this guide \\
\midrule
Sec.~1 & Related work on Kinect noise (Sec.~4.1) \\
Sec.~2.1 & Noise model, Eq.~1--2, Kinect numbers (Sec.~4.2) \\
Sec.~2.2 & Empirical validation, $1/8$ px, Algorithm~1 (Sec.~4.3) \\
Sec.~3.1--3.2 & Adaptive bilateral filter, $\sigma_d=kZ^2$ (Sec.~4.4) \\
Sec.~3.3 & Weighted volumetric scan merging, $w=1/Z^4$ (Sec.~4.4) \\
Sec.~3.4--3.5 & Plane extraction from disparity (Sec.~4.4) \\
Sec.~4 & Conclusion (Sec.~4.5) \\
\bottomrule
\end{tabularx}
\end{center}

\end{document}
